\chapter{Conclusion and Future Directions}\label{Conclusion}
\thispagestyle{empty}
\setlength{\parskip}{10pt plus 5pt minus 3pt}
\setlength{\parindent}{20pt}%

\section{Conclusions}
In the previous chapters we learnt the various features and the classifiers that have been used in the area of human activity recognition. We discussed the proposed approach for classifying human actions and discussed the methodology to get the feature representation for videos. Then we looked at the working of the HIMK based SVM classification for video data.\\
We used KTH dataset for experimentation and after looking at the results, we realized that the proposed model works better than the existing string kernel based model. 

\section{Future Work}
Looking at the results in Table \ref{comparison}, it can be easily deduced that the proposed system along with some modifications can be really helpful in the area of human action recognition.\\
We propose the following possible areas where work can be done:-
\begin{enumerate}
\item Use of motion features that capture the velocity of the action can be done.
\item Currently, there has been a shift to the deep learning methodologies for the area of videos. The deep learning based features can be explored for the task of human action recogntion.
\item We can also look at various modifications to the HIMK as proposed by Dileep to get better results.
\end{enumerate}

